\chapter{Модель и методы}
\label{chap:model-and-methods}

\section{Модель трейса}

В рамках настоящей работы трейс рассматривается как связанная совокупность
операций, относящихся к выполнению одного пользовательского запроса или другого
логически целостного действия. Такое представление соответствует общему подходу
OpenTelemetry, в котором трейс объединяет несколько спанов, принадлежащих одной
причинно-связанной цепочке исполнения~\cite{otelOverview,otelTraceApi}.

С практической точки зрения трейс в исследуемой библиотеке связывается прежде
всего с обработкой одного входящего HTTP-запроса. Если сервис принимает запрос,
затем вызывает один или несколько зависимых сервисов, а те продолжают выполнение
цепочки, все эти действия должны рассматриваться как части одного трейса. Это
позволяет анализировать распределенный сценарий как единое целое, а не как набор
локальных журналов и независимых временных отрезков.

Для дальнейшего описания удобно задать трейс как четверку
\((S, R, I, U)\), где \(S\) --- множество спанов, \(R\) --- отношение
родительско-дочерней связи между ними, \(I\) --- общий идентификатор трейса, а
\(U\) --- пользовательский или прикладной сценарий, который этот трейс
представляет. Такая запись не вводит дополнительный формат данных в библиотеку,
а только фиксирует смысловые элементы модели. Множество \(S\) описывает все
наблюдаемые операции, отношение \(R\) задает порядок и причинность вызовов,
идентификатор \(I\) связывает элементы в единый контекст, а сценарий \(U\)
показывает, зачем эти операции рассматриваются совместно.

Из этой модели следует важное ограничение: отдельный спан сам по себе не равен
трейсу. Спан описывает одну операцию, а трейс описывает связанную структуру
операций. Поэтому при проектировании библиотеки необходимо одновременно
поддерживать создание локальных спанов и перенос контекста между процессами. Без
переноса контекста множество спанов может быть создано корректно, но отношение
между ними будет потеряно на границе сервисов.

Для различения трейсов используется идентификатор трейса. Он должен сохраняться при
переходе между сервисами и включаться в контекст передачи. В модели библиотеки
идентификатор трейса не создается и не передается как отдельная прикладная
сущность; он присутствует в контексте трассировки и распространяется через
стандартный механизм W3C Trace Context~\cite{w3cTraceContext}. Тем самым модель
трейса опирается не на собственный формат идентификаторов, а на совместимый
межсистемный стандарт.

Трейс удобно представлять как ориентированную структуру вызовов. В простейшем
случае она близка к дереву: корневой серверный спан соответствует входящему
запросу, а дочерние спаны описывают вложенные шаги и исходящие обращения. Однако
для описания методов работы библиотеки достаточно использовать именно иерархическое
представление, поскольку оно отражает причинную связь между входящим запросом,
внутренними действиями сервиса и зависимыми вызовами.

\section{Модель спана}

Спан является базовой единицей наблюдаемого выполнения. Он описывает отдельную
операцию, имеющую начало, завершение, имя и набор метаданных~\cite{otelTraceApi,otelSdkSpan}.
В библиотеке такие операции соответствуют, например, приему входящего HTTP-запроса
или отправке исходящего HTTP-вызова.

С точки зрения логической модели спан содержит следующие группы данных:
\begin{enumerate}
    \item принадлежность одному трейсу;
    \item собственный идентификатор;
    \item ссылку на родительский спан, если операция является вложенной;
    \item временные границы выполнения;
    \item атрибуты операции;
    \item события, фиксирующие значимые внутренние моменты.
\end{enumerate}

В исследуемой реализации серверный спан создается в прослойке `MyOtelTracingLayer`
для каждого входящего HTTP-запроса. Он получает имя `http.request` и сразу
содержит поля, связанные с видом операции: метод запроса, путь, маршрут, код
ответа, длительность и признаки ошибки. Клиентский спан создается при отправке
исходящего запроса в `TracedRequestBuilder::send` и получает имя
`http.client.request`. Такой подход задает единообразную модель для двух главных
границ распределенного взаимодействия: входа запроса в сервис и выхода запроса к
зависимой системе.

На уровне модели важно не столько точное имя поля, сколько то, что библиотека
рассматривает спан как контейнер наблюдаемой операции. В него записываются
атрибуты HTTP-запроса, информация о результате выполнения и, при необходимости,
дополнительные бизнес-события. Следовательно, спан выступает не только как
временной интервал, но и как носитель контекста диагностики.

\begin{table}[H]
    \centering
    \caption{Поля спана и их роль в модели}
    \label{tab:span-fields}
    \renewcommand{\arraystretch}{1.2}
    \begin{tabularx}{\textwidth}{>{\raggedright\arraybackslash}p{3.2cm} >{\raggedright\arraybackslash}X >{\raggedright\arraybackslash}X}
        \toprule
        Поле модели & Назначение & Пример для HTTP-сервиса \\
        \midrule
        Имя операции & Кратко определяет наблюдаемое действие & `http.request`, `http.client.request` \\
        \addlinespace
        Временные границы & Позволяют вычислить длительность операции и порядок выполнения & момент приема запроса и момент формирования ответа \\
        \addlinespace
        Идентификатор спана & Отличает одну операцию от другой внутри одного трейса & идентификатор серверного или клиентского спана \\
        \addlinespace
        Родительский спан & Восстанавливает причинную связь между операциями & клиентский спан исходящего вызова является дочерним к серверному спану обработки \\
        \addlinespace
        Атрибуты & Описывают свойства операции в виде пар ключ--значение & метод, путь, маршрут, код ответа, удаленный узел \\
        \addlinespace
        События & Фиксируют значимые моменты внутри операции & начало бизнес-операции, завершение шага, доменное событие \\
        \addlinespace
        Статус & Показывает результат выполнения операции & успешное завершение или ошибка HTTP-класса 4xx/5xx \\
        \bottomrule
    \end{tabularx}
\end{table}

\begin{table}[H]
    \centering
    \caption{Основные сущности модели и их содержимое}
    \label{tab:model-entities}
    \renewcommand{\arraystretch}{1.2}
    \begin{tabularx}{\textwidth}{>{\raggedright\arraybackslash}p{2.7cm} >{\raggedright\arraybackslash}X >{\raggedright\arraybackslash}X}
        \toprule
        Сущность & Основные данные & Назначение \\
        \midrule
        Трейс & Идентификатор трейса, набор связанных спанов & Представление одного распределенного сценария выполнения \\
        \addlinespace
        Спан & Идентификатор спана, ссылка на родителя, временные границы, атрибуты, события & Описание отдельной операции внутри трейса \\
        \addlinespace
        Контекст спана & Идентификатор трейса, идентификатор спана, флаги трассировки, дополнительные сведения состояния & Передача минимального переносимого контекста между сервисами \\
        \bottomrule
    \end{tabularx}
\end{table}

\section{Контекст спана}

Контекст спана представляет собой минимальный набор данных, который должен быть
передан между независимыми процессами для продолжения уже начатого трейса. В
соответствии со стандартом W3C Trace Context и средствами OpenTelemetry такой
контекст включает идентификатор трейса, идентификатор текущего спана и флаги
трассировки~\cite{w3cTraceContext,opentelemetryPropagationDocs}.

В библиотеке чтение и запись контекста выполняются через модуль `propagation`.
Функция `extract_context` извлекает контекст из набора HTTP-заголовков, а функция
`inject_context` записывает текущий контекст в заголовки исходящего запроса. В
качестве адаптера используются структуры `HeaderExtractor` и `HeaderInjector`,
которые связывают абстракции OpenTelemetry с типом `http::HeaderMap`.

Такой подход важен по двум причинам. Во-первых, библиотека не реализует
собственный формат распространения контекста, а использует стандартный
распространитель текста OpenTelemetry. Во-вторых, работа с контекстом изолирована в
отдельном модуле, поэтому правила переноса трейса не размазываются по коду
прослойки и клиентских вызовов. Это упрощает верификацию поведения и снижает риск
рассогласования между входящей и исходящей обработкой.

Переносимый контекст должен быть компактным и технически нейтральным. Он должен
сериализоваться в HTTP-заголовки, поскольку именно заголовки являются общей
границей взаимодействия для сервисов, шлюзов и промежуточных прокси. При этом
контекст не должен превращаться в контейнер прикладных данных. Идентификаторы
трассировки нужны для связывания операций, а не для передачи сведений о
пользователе, заказе или состоянии бизнес-процесса.

Отдельное ограничение связано с безопасностью. В переносимый контекст не должны
попадать чувствительные значения: токены доступа, файлы cookie, персональные
идентификаторы и другие данные, которые не требуются для восстановления
родительско-дочерней связи. Такие сведения, если они действительно нужны для
диагностики, должны обрабатываться отдельной политикой захвата атрибутов:
разрешительным списком, усечением, маскированием или фиксацией только факта
наличия.

\section{Родительско-дочерние связи}

Для корректного восстановления структуры распределенного сценария необходимо явно
задавать связь между родительскими и дочерними спанами. В модели библиотеки
используются три характерных роли операций:
\begin{enumerate}
    \item корневая или продолжаемая серверная операция при обработке входящего запроса;
    \item вложенные операции внутри сервиса;
    \item клиентские операции при обращении к зависимым HTTP-сервисам.
\end{enumerate}

Если во входящих заголовках уже присутствует контекст трассировки, серверный спан
должен продолжать существующий трейс. Если контекст отсутствует, создается новый
трейс, в котором входящий серверный спан становится фактической точкой начала
наблюдаемого сценария. Далее все вложенные действия внутри обработчика выполняются
в пределах этого спана или его дочерних элементов.

При формировании исходящего HTTP-запроса клиентский спан должен быть логически
дочерним по отношению к активной операции текущего сервиса. После этого именно его
контекст записывается в заголовки запроса, чтобы зависимый сервис мог продолжить
общий трейс. Благодаря этому входящий запрос в одном сервисе и исходящий запрос к
другому сервису оказываются связанными в одной иерархии.

\begin{figure}[H]
    \centering
    \begin{tikzpicture}[
        node distance=1.9cm,
        block/.style={draw, rounded corners, minimum width=3.5cm, minimum height=0.9cm, align=center},
        arrow/.style={->, thick}
    ]
        \node[block] (trace) {Трейс};
        \node[block, below left=of trace] (server) {Серверный спан\\входящего запроса};
        \node[block, below right=of trace] (client) {Клиентский спан\\исходящего запроса};
        \node[block, below=of server] (event) {События и атрибуты\\внутри обработки};

        \draw[arrow] (trace) -- (server);
        \draw[arrow] (trace) -- (client);
        \draw[arrow] (server) -- (event);
        \draw[arrow] (server) -- (client);
    \end{tikzpicture}
    \caption{Логическая связь трейса, серверного спана, клиентского спана и внутреннего контекста обработки}
    \label{fig:trace-model}
\end{figure}

\section{Метод обработки входящего HTTP-запроса}

Метод обработки входящего запроса в библиотеке строится вокруг прослойки
`MyOtelTracingLayer`, которая оборачивает внутренний сервис. С точки зрения
алгоритма этот процесс можно представить следующей последовательностью:
\begin{enumerate}
    \item принять HTTP-запрос на границе сервиса;
    \item определить метод и путь запроса, а при наличии данных маршрутизации ---
    шаблон маршрута;
    \item извлечь контекст трассировки из заголовков запроса;
    \item создать серверный спан для операции `http.request`;
    \item привязать к спану маршрут, метод, путь и разрешенные атрибуты заголовков;
    \item выполнить внутренний обработчик внутри созданного спана;
    \item после завершения обработки записать код ответа, длительность и признаки ошибки.
\end{enumerate}

Реализация в `layer.rs` подтверждает эту модель. Прослойка фиксирует момент начала
обработки, создает спан `http.request`, затем, если трассировка через
OpenTelemetry включена, устанавливает родительский контекст из входящих заголовков.
После этого внутренний обработчик выполняется через `instrument(span.clone())`,
что обеспечивает наследование контекста для вложенных действий. Когда обработка
завершается, в спан записываются код ответа и длительность, а при клиентской или
серверной ошибке --- дополнительные признаки ошибочного завершения.

С методической точки зрения такой подход важен тем, что логика наблюдаемости
расположена на границе обработки, а не в каждом отдельном обработчике. Это
обеспечивает единообразие поведения для всех маршрутов сервиса и уменьшает объем
повторяющегося кода.

Псевдокод метода обработки входящего запроса можно записать следующим образом:

\begin{verbatim}
receive request
extract parent context from request headers
create server span "http.request"
record method, path, route and allowed request attributes
run inner service inside server span
record response status, duration and error attributes
return response
\end{verbatim}

Этот алгоритм намеренно описан на уровне действий, а не на уровне конкретных
операторов Rust. Его задача состоит в том, чтобы отделить метод обработки от
деталей реализации. В главе 4 этот же порядок действий будет сопоставлен с
модулем `layer.rs` и фактическим кодом прослойки.

\section{Метод обработки исходящего HTTP-запроса}

Исходящий HTTP-запрос в библиотеке рассматривается как отдельная наблюдаемая
операция, которая должна быть связана с текущим активным трейсом. Для этого
используется специализированный клиентский тип, возвращающий построитель
трассируемого запроса.

Логика метода отправки запроса включает следующие шаги:
\begin{enumerate}
    \item сформировать исходящий HTTP-запрос на основе метода, адреса и тела;
    \item создать клиентский спан `http.client.request`;
    \item записать в него сведения о методе, пути и адресе удаленного узла;
    \item внедрить текущий контекст трассировки в заголовки запроса;
    \item отправить запрос;
    \item после получения ответа записать код ответа, длительность и признаки ошибки;
    \item в случае сетевого сбоя зафиксировать ошибочный тип завершения.
\end{enumerate}

В `request_builder.rs` эта схема реализована последовательно. После разбиения
построителя на клиент и готовый запрос библиотека создает клиентский спан,
внедряет контекст через `inject_context`, а затем выполняет запрос внутри этого
спана. При успешном ответе записываются код и длительность, а при ошибке
устанавливается тип ошибки `reqwest`. За счет этого внешние вызовы становятся
видимыми как части общего трейса и могут быть сопоставлены с соответствующими
входящими операциями зависимых сервисов.

Псевдокод исходящего запроса имеет следующий вид:

\begin{verbatim}
build outgoing request
create client span "http.client.request"
record method, path and remote host
inject current trace context into request headers
send request inside client span
record response status, duration and error attributes
return response or typed error
\end{verbatim}

В отличие от входящего алгоритма, здесь ключевым шагом является внедрение
контекста перед отправкой запроса. Если этот шаг пропустить, зависимый сервис
сможет создать собственный серверный спан, но он не будет связан с исходным
пользовательским сценарием. Поэтому метод исходящего запроса должен рассматриваться
как обязательная часть модели распределенной трассировки, а не как вспомогательная
обертка над HTTP-клиентом.

\section{Передача контекста, атрибуты и события}

Передача контекста, запись атрибутов и фиксация событий образуют дополнительные
методы обогащения трейса, без которых распределенная трассировка теряет
диагностическую ценность.

Передача контекста реализуется через стандарт W3C Trace Context и модуль `propagation`.
Это гарантирует, что библиотека переносит не произвольный набор полей, а
стандартизованный контекст, который может быть интерпретирован другими
совместимыми системами~\cite{w3cTraceContext,opentelemetryPropagationDocs}.

Запись атрибутов в библиотеке разделена на две группы. Первая группа формируется
автоматически и относится к HTTP-операции: метод, путь, маршрут, код ответа,
адрес удаленного узла, длительность и признаки ошибки. Вторая группа определяется
пользователем через правила захвата заголовков. Для этого используется
`HeaderCapturePolicy`, которая задает:
\begin{enumerate}
    \item список разрешенных заголовков;
    \item правило преобразования имени заголовка в ключ атрибута;
    \item режим записи значения;
    \item максимальную длину значения;
    \item максимальное число захватываемых заголовков;
    \item поведение для не-UTF-8 значений.
\end{enumerate}

Поддерживаются четыре режима записи значения заголовка: полное сохранение,
усечение до заданной длины, маскирование и фиксация одного факта присутствия.
Дополнительно существует список заголовков, которые запрещено записывать по
соображениям безопасности, например `authorization`, `cookie` и `x-api-key`.
Тем самым модель библиотеки предполагает, что захват атрибутов является
разрешительным, а не всеобъемлющим.

События фиксируются через функцию `record_event`. Она принимает имя события и набор
полей, преобразует их в структурированное представление и записывает как
информационное сообщение трассировки с именем события и сериализованными
значениями. Доступные типы значений включают строку, логическое значение, целое и
вещественное число. Это позволяет встраивать в трейс прикладной контекст без
ручного конструирования низкоуровневых записей.

\section{Конфигурация построителя}

Центральной конфигурационной сущностью библиотеки является `TracingConfig`, которая
создается через `TracingConfig::builder`. Этот подход задает метод поэтапного
построения конфигурации и отделяет настройку параметров от их последующей
валидации.

На основании реализации `config.rs` можно выделить следующие основные параметры:
\begin{enumerate}
    \item имя сервиса;
    \item версия сервиса;
    \item окружение выполнения;
    \item адрес конечной точки OTLP;
    \item фильтр журналирования;
    \item дополнительные атрибуты ресурса;
    \item режим выборки трейсов;
    \item тайм-аут экспорта;
    \item тайм-аут завершения работы;
    \item формат журналов.
\end{enumerate}

По умолчанию имя сервиса должно быть задано явно, окружение принимает значение
`local`, конечная точка OTLP указывает на локальный приемник по адресу
`http://localhost:4318/v1/traces`, режим выборки всегда включен, а формат журналов
по умолчанию --- структурированный. При построении конфигурации выполняется валидация:
проверяется непустое имя сервиса, корректность адреса конечной точки OTLP,
допустимость пользовательских ключей атрибутов и ненулевые значения тайм-аутов.

Таким образом, построитель выполняет не только функцию удобной настройки, но и
роль защитного барьера, который не допускает запуск библиотеки с заведомо
некорректной конфигурацией.

Ошибки конфигурации в модели делятся на несколько групп. Первая группа относится к
обязательным значениям: имя сервиса не может быть пустым, поскольку оно
используется для идентификации источника телеметрии. Вторая группа относится к
адресам и временным ограничениям: конечная точка OTLP должна иметь допустимую
схему и узел назначения, а тайм-ауты экспорта и завершения не могут быть
нулевыми. Третья группа связана с пользовательскими атрибутами: ключи не должны
нарушать правила именования и не должны пересекаться с зарезервированными
служебными полями.

Такое разделение важно для пользователя библиотеки. Ошибка должна возникать до
запуска подсистемы трассировки, а не после первого запроса или при завершении
процесса. Поэтому метод построения конфигурации завершает не только создание
структуры `TracingConfig`, но и проверку минимальных условий корректной работы.
Если условия не выполнены, прикладной сервис получает типизированную ошибку и
может прекратить запуск с понятной причиной.

\section{Ограничения минимально необходимой версии}

Минимально необходимая версия библиотеки ориентирована на решение конкретной
практической задачи: подключение распределенной трассировки к HTTP-сервисам на
Rust с минимальным объемом повторяющегося кода. Поэтому ее модель сознательно
ограничена рядом условий.

\begin{table}[H]
    \centering
    \caption{Ограничения минимально необходимой версии библиотеки}
    \label{tab:mvp-limitations}
    \renewcommand{\arraystretch}{1.2}
    \begin{tabularx}{\textwidth}{>{\raggedright\arraybackslash}p{4.4cm} >{\raggedright\arraybackslash}X}
        \toprule
        Ограничение & Смысл ограничения \\
        \midrule
        Ориентация на HTTP & Рассматриваются входящие и исходящие HTTP-вызовы, а не произвольные транспортные протоколы \\
        \addlinespace
        Основной стек Axum и Tower & Прослойка строится вокруг существующих абстракций этих библиотек, а не вокруг всех возможных каркасов \\
        \addlinespace
        Экспорт через OTLP/HTTP & Экспорт ориентирован на совместимую инфраструктуру наблюдаемости и не включает собственное хранилище трейсов \\
        \addlinespace
        Ограниченный режим выборки & На текущем этапе реализован только режим постоянной выборки \\
        \addlinespace
        Частичное покрытие телеметрии & Основной акцент сделан на трейсах и корреляции журналов, а не на полном наборе метрик и журналов как самостоятельных сигналов \\
        \bottomrule
    \end{tabularx}
\end{table}

Эти ограничения не противоречат цели работы, поскольку она состоит не в создании
универсальной платформы наблюдаемости, а в разработке библиотечного слоя для
низкопорогового подключения распределенной трассировки в типичном HTTP-сценарии
на Rust. Напротив, ограничение области применения позволяет сделать модель
библиотеки достаточно определенной и пригодной для последующей реализации.

Отсутствие метрик в минимально необходимой версии не нарушает основную цель,
поскольку задача работы состоит в восстановлении пути запроса между сервисами.
Метрики полезны для агрегированного мониторинга, но сами по себе не показывают
структуру конкретного пользовательского сценария. По той же причине журналы в
данной версии рассматриваются прежде всего как канал корреляции с активным
трейсом, а не как отдельный экспортируемый сигнал OpenTelemetry.

Не включается также перенос baggage и поддержка всех возможных веб-каркасов.
Baggage предназначен для распространения пользовательских пар ключ--значение и
требует отдельной политики безопасности. Поддержка других каркасов, например
Actix Web или gRPC-стека, расширила бы область проектирования, но не изменила бы
базовую модель входящего и исходящего HTTP-запроса. Поэтому минимальная версия
сосредоточена на Axum, Tower, OTLP/HTTP и Jaeger как достаточном контуре для
проверки выбранного подхода.

В результате по данной главе можно сформулировать два методических вывода. Во-первых,
библиотека должна опираться на четкую модель трейса, спана и переносимого
контекста. Во-вторых, методы обработки входящих и исходящих HTTP-запросов должны
быть организованы так, чтобы передача контекста, запись атрибутов и фиксация
ошибок выполнялись автоматически и единообразно. Эти выводы служат основанием для
следующей главы, в которой рассматривается проектирование архитектуры библиотеки.
